% Computer Science A --- Spring Framework (Core, MVC, Data). Boot-specific items live in \texttt{springboot.tex}.

\runningheader{Computer Science A}{Spring Framework}{Page \thepage\ of \numpages}

\begingradingrange{csaspring}

\question[4]
Spring is
\begin{choices}
  \choice A dependency injection framework for JavaScript.
  \choice A mocking library for Java.
  \choice A JVM based programming language.
  \CorrectChoice A dependency injection framework for Java.
\end{choices}
\begin{solution}
  The \textbf{Spring Framework} for Java provides \textbf{dependency injection}, declarative middleware (transactions, AOP), and modules such as Spring MVC and Spring Data. It is not a language, not JavaScript-focused, and not primarily a mocking library.
\end{solution}

\question[4]
A Spring bean is
\begin{choices}
  \choice The same as a Java bean.
  \CorrectChoice A class which is constructed and managed by Spring.
  \choice Only classes defined in the Spring source code.
  \choice Any bean annotated with \texttt{@GetMapping}.
\end{choices}
\begin{solution}
  A \textbf{bean} is an object whose lifecycle (construction, wiring, destruction) is handled by the Spring \textbf{IoC container}. JavaBeans are a naming/convention pattern; \texttt{@GetMapping} maps HTTP, not bean identity.
\end{solution}

\question[4]
Which of the following is \textbf{not} a way to get a bean from Spring?
\begin{choices}
  \choice Using the \texttt{getBean} method on the \texttt{ApplicationContext} or \texttt{BeanFactory}.
  \choice Using \texttt{@Autowired} with a valid bean definition.
  \choice Using an XML configuration with a valid bean reference.
  \CorrectChoice Using the \texttt{@Aspect} annotation before the declaration field, setter, or constructor.
\end{choices}
\begin{solution}
  Lookup uses the context (\texttt{getBean}), injection (\texttt{@Autowired} / constructor injection), or XML references. \texttt{@Aspect} declares AOP aspects, not bean retrieval. (Source used \texttt{@Autowire}; standard Spring spelling is \texttt{@Autowired}.)
\end{solution}

\question[4]
Which of the following is \textbf{not} a valid stereotype annotation?
\begin{choices}
  \choice \texttt{@Component}
  \CorrectChoice \texttt{@Bean}
  \choice \texttt{@Controller}
  \choice \texttt{@Service}
\end{choices}
\begin{solution}
  \textbf{Stereotypes} (\texttt{@Component}, \texttt{@Service}, \texttt{@Repository}, \texttt{@Controller}, \ldots) mark classes for component scanning. \texttt{@Bean} is used on \textbf{factory methods} (often inside \texttt{@Configuration}) to register a bean, not as a class-level stereotype.
\end{solution}

\question[4]
When requesting a bean with the scope \texttt{prototype}, you will always be given a new instance of the bean.
\begin{choices}
  \CorrectChoice true
  \choice false
\end{choices}
\begin{solution}
  \textbf{Prototype} scope returns a \textbf{new} instance for each \texttt{getBean} or injection point (depending on how it is resolved), unlike \textbf{singleton}, where one shared instance is typical.
\end{solution}

\question[4]
The default scope of a bean is \texttt{singleton}.
\begin{choices}
  \CorrectChoice true
  \choice false
\end{choices}
\begin{solution}
  In Spring, the default scope for beans in the IoC container is \textbf{singleton}: one logical instance per container (for that bean name).
\end{solution}

\question[4]
One of the problems Spring Core addresses is
\begin{choices}
  \CorrectChoice Tight coupling
  \choice Loose coupling
  \choice Pagination
  \choice Session management
\end{choices}
\begin{solution}
  Spring's \textbf{Inversion of Control} and \textbf{dependency injection} reduce \textbf{tight coupling} by injecting collaborators instead of hard-coding \texttt{new}. The goal is \emph{loose} coupling; pagination and sessions are separate concerns.
\end{solution}

\question[4]
The \texttt{@Configuration} annotation is used to
\begin{choices}
  \choice Provide XML configuration.
  \CorrectChoice Define a class as a configuration bean.
  \choice Provide annotation driven configuration.
  \choice Do the same as the \texttt{@Autowired} annotation.
\end{choices}
\begin{solution}
  \texttt{@Configuration} marks a class as a source of \textbf{@Bean} definitions (Java-based config). It does not replace all XML by itself; \texttt{@Autowired} declares dependency injection points, not configuration structure.
\end{solution}

\question[4]
You can provide configuration to Spring using XML, a configuration bean, and annotation-driven configuration.
\begin{choices}
  \CorrectChoice true
  \choice false
\end{choices}
\begin{solution}
  Spring supports \textbf{XML}, \textbf{Java config} (\texttt{@Configuration} classes), and \textbf{annotation-driven} setup (e.g.\ component scanning with \texttt{@Component}, \texttt{@Service}, etc.), often combined in one application.
\end{solution}

\question[4]
Which of the following dependency injection types are supported by the Spring Framework?
\begin{choices}
  \choice setter injection
  \choice constructor injection
  \CorrectChoice Both of the above.
  \choice None of the above.
\end{choices}
\begin{solution}
  Spring supports \textbf{constructor} and \textbf{setter} injection (and field injection via \texttt{@Autowired}). Interface injection is less central in modern Spring.
\end{solution}

\question[4]
Which of the following is an advantage of using Setter Injection in Spring?
\begin{choices}
  \choice It ensures that all dependencies are provided at the time of object creation.
  \CorrectChoice It allows for optional dependencies to be injected.
  \choice It allows for dependencies to be set using reflection.
  \choice None of the above
\end{choices}
\begin{solution}
  Setters allow \textbf{optional} collaborators to be omitted or changed after construction. \textbf{Constructor} injection is preferred for required, immutable dependencies.
\end{solution}

\question[4]
The \texttt{@RequestMapping} annotation is used to
\begin{choices}
  \CorrectChoice Provide a URL mapping to a resource for any HTTP method.
  \choice Provide a URL mapping to a resource for only GET methods.
  \choice Map a path variable to an array index.
  \choice Provide data contained in a request body as a \texttt{HashMap}.
\end{choices}
\begin{solution}
  \texttt{@RequestMapping} can map a path (and optionally HTTP methods via \texttt{method = \ldots}). Narrower annotations (\texttt{@GetMapping}, \texttt{@PostMapping}, \ldots) are specializations. Path variables use \texttt{@PathVariable}; the body is bound with \texttt{@RequestBody} or parameters, not \texttt{@RequestMapping} alone.
\end{solution}

\question[4]
What is the \texttt{@Controller} annotation?
\begin{choices}
  \CorrectChoice The \texttt{@Controller} annotation indicates that a particular class serves the role of a controller.
  \choice The \texttt{@Controller} annotation indicates how to control the transaction management.
  \choice The \texttt{@Controller} annotation indicates how to control the dependency injection.
  \choice The \texttt{@Controller} annotation indicates how to control the aspect programming.
\end{choices}
\begin{solution}
  In Spring MVC, \texttt{@Controller} (stereotype of \texttt{@Component}) marks a web controller whose methods handle HTTP via \texttt{@RequestMapping} and related annotations. It does not govern transactions, DI, or AOP by itself.
\end{solution}

\question[4]
When used on the class, the \texttt{@RestController} annotation implies what on all methods?
\begin{choices}
  \CorrectChoice \texttt{@ResponseBody}
  \choice \texttt{@RequestBody}
  \choice \texttt{@RequestMapping}
  \choice \texttt{@RequestParam}
\end{choices}
\begin{solution}
  \texttt{@RestController} is \texttt{@Controller} + \texttt{@ResponseBody}: return values of handler methods are serialized to the HTTP response body (e.g.\ JSON) by default. \texttt{@RequestBody} binds incoming body; \texttt{@RequestMapping} maps URLs; \texttt{@RequestParam} binds query/form params.
\end{solution}

\question[4]
The \texttt{@RequestBody} annotation is used to
\begin{choices}
  \choice Define a class as a controller.
  \choice Define that the return value of a method should be serialized to some form such as XML, JSON, ect.
  \CorrectChoice Define that the request body should be deserialized to a Java object from some form such as XML, JSON, ect.
  \choice Parse a variable in the requested URL and inject it into the method call as a parameter.
\end{choices}
\begin{solution}
  \texttt{@RequestBody} binds the \textbf{HTTP request body} to a Java object using message converters (often JSON). Response serialization uses \texttt{@ResponseBody} or \texttt{@RestController}. (Original uses \textit{ect.})
\end{solution}

\question[4]
The \texttt{@PathVariable} annotation is used to
\begin{choices}
  \choice Define a class as a controller.
  \choice Define that the return value of a method should be serialized to some form such as XML, JSON, ect.
  \choice Define that the request body should be deserialized to a Java object from form such as XML, JSON, ect.
  \CorrectChoice Parse the variable in the requested URL and provide it to the method call as a parameter.
\end{choices}
\begin{solution}
  \texttt{@PathVariable} binds URI template variables (e.g.\ \texttt{/orders/\{id\}}) to method parameters. It is not for the request body or return serialization.
\end{solution}

\question[4]
Spring MVC. Given the following annotation, what would be the most appropriate handler syntax?
\texttt{@GetMapping("/\{id\}")}
\begin{choices}
  \choice \texttt{public Object getById(@PathParam(name="id") int value) \{\}}
  \choice \texttt{public Object getById(@RequestBody int id) \{\}};
  \choice \texttt{public Object getById(@RequestMapping int id) \{\}}
  \CorrectChoice \texttt{public Object getById(@PathVariable int id) \{\}}
\end{choices}
\begin{solution}
  Path segments like \texttt{\{id\}} bind with Spring's \texttt{@PathVariable}. \texttt{@PathParam} is JAX-RS, not Spring MVC. \texttt{@RequestBody} is for the HTTP body; \texttt{@RequestMapping} is not a parameter annotation.
\end{solution}

\question[4]
The \texttt{@Transactional} annotation is used to
\begin{choices}
  \CorrectChoice Define that the underlying code should be executed as a database transaction.
  \choice Define that the underlying code is used to make a purchase.
  \choice Define that the underlying code should never throw an exception.
  \choice Does nothing.
\end{choices}
\begin{solution}
  \texttt{@Transactional} (Spring's declarative transaction support) demarcates a \textbf{transaction} boundary, typically around service methods hitting a transactional resource. Behavior on exceptions depends on rollback rules.
\end{solution}

\question[4]
Spring Data's \texttt{JpaRepository} features a fluent API which allows it to
\begin{choices}
  \choice Generate Procedural Langauge functions on the database.
  \choice Know whether data exists without checking the database.
  \choice Provide an in-memory database that doesn't require a separate database.
  \CorrectChoice Generate SQL queries by parsing an abstract method's name.
\end{choices}
\begin{solution}
  \textbf{Query methods} (derived from method names like \texttt{findByLastName}) let Spring Data JPA build queries from signatures. The first choice keeps the original typo \textit{Langauge}. Existence still hits the database (e.g.\ \texttt{existsBy\ldots}); in-memory DB is not what \texttt{JpaRepository} itself provides.
\end{solution}

\question[4]
Regarding Spring Data, what are custom query methods?
\begin{choices}
  \choice Methods whose names define what query to perform using property expressions. You will need to create a class that implements these methods.
  \CorrectChoice Methods whose names define what query to perform using property expressions. Spring will automatically generate the implementation for you.
  \choice Methods that allow developers to write SQL queries directly in Java code.
  \choice Predefined methods provided by Spring Data for common database operations.
\end{choices}
\begin{solution}
  \textbf{Derived query methods} on a repository interface are implemented by the Spring Data infrastructure at runtime from the method name and domain model. You can also use \texttt{@Query} for explicit JPQL/SQL.
\end{solution}

\uplevel{\par\textbf{Spring Data annotations} (from the fill-in table in the source bank).\par}

\question[4]
Which annotation marks a Spring Data repository \emph{interface} that should \textbf{not} be picked up as a Spring Data repository implementation (for example a shared base interface with only common signatures)?
\begin{choices}
  \CorrectChoice \texttt{@NoRepositoryBean}
  \choice \texttt{@Repository}
  \choice \texttt{@Query}
  \choice \texttt{@Id}
\end{choices}
\begin{solution}
  \texttt{@NoRepositoryBean} tells Spring Data not to create a proxy for that interface---useful for intermediate templates. Concrete repository interfaces (extending \texttt{JpaRepository}, etc.) are normally discovered without this marker.
\end{solution}

\question[4]
Which annotation marks a field in a JPA entity class as the primary key? (The source table had a typo \texttt{@ld}.)
\begin{choices}
  \choice \texttt{@NoRepositoryBean}
  \choice \texttt{@ld}
  \CorrectChoice \texttt{@Id}
  \choice \texttt{@Transient}
\end{choices}
\begin{solution}
  JPA uses \texttt{@Id} on the primary-key field or property. \texttt{@ld} is a typo in the original key column; \texttt{@Transient} means not persisted.
\end{solution}

\question[4]
Which annotation marks a field to be ignored by the persistence provider during reads and writes (not stored in the database)?
\begin{choices}
  \choice \texttt{@Id}
  \choice \texttt{@Query}
  \CorrectChoice \texttt{@Transient}
  \choice \texttt{@NoRepositoryBean}
\end{choices}
\begin{solution}
  \texttt{javax.persistence.Transient} (or \texttt{jakarta.persistence.Transient}) excludes the field from the entity mapping. Spring Data MongoDB also has field-level transient support in its mapping model.
\end{solution}

\question[4]
Which annotation supplies a JPQL (or native) query string for a Spring Data repository method?
\begin{choices}
  \choice \texttt{@NoRepositoryBean}
  \choice \texttt{@Id}
  \choice \texttt{@Transient}
  \CorrectChoice \texttt{@Query}
\end{choices}
\begin{solution}
  \texttt{@Query} on a repository method declares JPQL/SQL explicitly when derived names are insufficient. Parameters can be bound by position or \texttt{@Param}.
\end{solution}

\question[4]
To use the Spring framework you must use Spring Boot
\begin{choices}
  \choice TRUE
  \CorrectChoice FALSE
\end{choices}
\begin{solution}
  The \textbf{Spring Framework} predates \textbf{Spring Boot}. Boot is a convention-over-configuration layer on top; plain Spring with XML or Java config remains valid without Boot.
\end{solution}

\question[4]
What is the purpose of using \texttt{MockMvc} for Spring MVC testing?
\begin{choices}
  \choice It is a replacement for JUnit and Mockito
  \CorrectChoice It can run integration tests that test controllers without explicitly starting a Servlet container
  \choice It can run regression tests that detect defects introduced by new code
  \choice It enables making calls to external application's REST API endpoints for testing
\end{choices}
\begin{solution}
  \texttt{MockMvc} (from \texttt{spring-test}) dispatches through the \textbf{DispatcherServlet} test infrastructure, so you exercise controller mapping, validation, and serialization without binding to a live port. It is often used from Spring Boot tests but is not a Boot-only API. JUnit/Mockito remain separate tools.
\end{solution}

\uplevel{\par\textbf{Spring Framework practice bank}\par}

\question[4]
Which statement best describes constructor injection in Spring?
\begin{choices}
  \CorrectChoice It makes required dependencies explicit at object creation and is generally preferred for mandatory collaborators
  \choice It is only available when XML configuration is used
  \choice It can only inject primitive values, not beans
  \choice It is deprecated in favor of field injection
\end{choices}
\begin{solution}
  Constructor injection works well for required dependencies and supports immutability/testing. Field injection still exists but is less explicit and harder to test in isolation.
\end{solution}

\question[4]
What is the most accurate difference between \texttt{@Controller} and \texttt{@RestController}?
\begin{choices}
  \choice \texttt{@Controller} cannot map HTTP requests
  \CorrectChoice \texttt{@RestController} implies \texttt{@ResponseBody} semantics on handler methods, while \texttt{@Controller} is typically used for MVC view rendering unless methods are annotated otherwise
  \choice \texttt{@RestController} is only valid in Spring Boot, not Spring Framework
  \choice They are aliases with no behavioral differences
\end{choices}
\begin{solution}
  \texttt{@RestController} combines \texttt{@Controller} and \texttt{@ResponseBody}, so return values are serialized to the response body by default.
\end{solution}

\question[4]
Given \texttt{@GetMapping("/orders/\{id\}")}, which parameter binding is correct?
\begin{choices}
  \choice \texttt{@RequestBody Long id}
  \CorrectChoice \texttt{@PathVariable Long id}
  \choice \texttt{@RequestParam Long id}
  \choice \texttt{@Autowired Long id}
\end{choices}
\begin{solution}
  URI template tokens (\texttt{\{id\}}) are bound with \texttt{@PathVariable}; \texttt{@RequestParam} binds query string/form parameters.
\end{solution}

\question[4]
Which statement about \texttt{@Transactional} default rollback behavior is correct?
\begin{choices}
  \CorrectChoice By default, Spring rolls back on unchecked exceptions (\texttt{RuntimeException} and \texttt{Error}) and commits on checked exceptions unless configured otherwise
  \choice It always rolls back for every thrown exception type
  \choice It only applies when methods are static
  \choice It opens one global JVM-wide transaction shared by all beans
\end{choices}
\begin{solution}
  You can customize rollback rules (e.g.\ \texttt{rollbackFor = Exception.class}) when checked exception rollback is needed.
\end{solution}

\question[4]
Which Spring Data repository method name will derive a query for users by exact email?
\begin{choices}
  \choice \texttt{queryUserEmail(String email)}
  \CorrectChoice \texttt{findByEmail(String email)}
  \choice \texttt{selectEmailFromUser(String email)}
  \choice \texttt{getEmailWhereUser(String email)}
\end{choices}
\begin{solution}
  Spring Data parses supported keywords/patterns (like \texttt{findBy...}, \texttt{existsBy...}) to generate query implementations.
\end{solution}

\question[4]
Which annotation should you use to inject a query parameter like \texttt{?page=2} into a controller method argument?
\begin{choices}
  \CorrectChoice \texttt{@RequestParam}
  \choice \texttt{@PathVariable}
  \choice \texttt{@RequestBody}
  \choice \texttt{@Configuration}
\end{choices}
\begin{solution}
  Query string values are bound via \texttt{@RequestParam}; body payloads use \texttt{@RequestBody}; path segments use \texttt{@PathVariable}.
\end{solution}

\endgradingrange{csaspring}
